\section{Methods}\label{sec:methods}

\subsection{Comparing cars within a lap}\label{sec:within}

Let $y_{rt}$ be the lap time on race lap $t$ of run $r$, driven by driver $j$ at event $e$ on compound $c(r)$ with tyre age $a_{rt}$ and traffic $q_{rt}$, and let $D_{c,r}$ indicate that run $r$ uses compound $c$.
In a race, fuel mass, track evolution and weather are nearly the same for every car on a given lap, so we write their joint contribution as a term $\lambda_{et}$ shared by all cars in the (event, lap) group.
For any quantity $z$, let $\tilde z_{rt}$ be $z$ minus its mean over the cars in the same group.
Applying this to every term of the lap-time model removes $\lambda_{et}$ without modelling it:
\begin{equation}\label{eq:race}
  \tilde y_{rt} = \sum_{c} \mu_c\,\tilde D_{c,r} + \sum_{c} \beta_{c}\,\widetilde{(D_c a)}_{rt} + \gamma\,\tilde q_{rt}
  + b_{e}\,\tilde a_{rt} + \tilde d_{j(r)e,t} + \tilde g_{r,t} + \widetilde{h_r a}_{rt} + \tilde\varepsilon_{rt}.
\end{equation}
Here $\beta_c$ is the degradation rate of compound $c$ in seconds per lap of tyre age and $\mu_c$ its pace offset, with one compound as the reference.
Every column is demeaned, including the compound indicators and their products with tyre age.
A tempting shortcut, multiplying each indicator by tyre age demeaned over all cars, is exact only when every car on the lap uses the same compound; when compounds are mixed it leaks one compound's slope and offset into another's, and on our data it moves the 2026 C1 rate from \RateCOne\spl{} to \ApproxSixRateOne\spl{} (Appendix~\ref{app:exact}).
The random effects are a circuit slope $b_e\sim N(0,\sigma_b^2)$, a driver-in-event intercept $d_{je}$, a run intercept $g_r$, a run slope $h_r$ and lap-level noise $\varepsilon_{rt}$, each independent and normal; tildes on these terms mean the same group demeaning of the car-level quantity.
Tyre age enters linearly.

The identifying assumption is that, within a lap, differences in tyre age between cars are not correlated with unmodelled pace once run intercepts and traffic are included.
Fuel is common to the group, so its coefficient is not needed and cannot be mis-specified.
Run intercepts matter because the cars on older tyres at a given lap are not a random draw: strategy and car pace decide who has stopped.
That selection shifts a car's pace level, which the run intercept absorbs, so the slope comes from how a run's lap time changes with tyre age relative to the field on the same laps.
A cross-sectional alternative that compares cars of different tyre age within a lap and compound without run terms picks up the selection and returns negative soft slopes (Appendix~\ref{app:xs}).
Demeaning makes the errors of cars in one group correlated, because they share the subtracted mean; we do not model this correlation, but intervals that allow arbitrary correlation within each event are essentially no wider than ours (Appendix~\ref{app:inference}).
Traffic is also affected by tyre age, since a car on worn tyres is more likely to be caught, so its coefficient is not a clean control (Section~\ref{sec:limits}).

Practice sessions offer no common lap, since cars run at different times.
Fuel and track evolution then both fall through the session, and the fuel coefficient is not identified from lap times.
We therefore subtract an assumed fuel effect and centre within the run:
\begin{equation}\label{eq:practice}
  z_{rt} = y_{rt} + \phi\,\kappa_{rt}, \qquad \dot z_{rt} = z_{rt}-\bar z_{r},
\end{equation}
where $\kappa_{rt}$ is the fuel burned since the run began, taken as 1.7\,kg per lap, and $\phi=\FuelCoefficient$\,s/kg.
These two constants are assumptions that we could not tie to an open source; Appendix~\ref{app:sens} varies both.
We regress $\dot z_{rt}$ on tyre age centred within the run.
The same assumed correction is added to every slope, so it cancels in differences between contexts and between compounds but not in ratios.

\subsection{Estimation and inference}\label{sec:mixed}

We fit Eq.~\eqref{eq:race} by restricted maximum likelihood (REML) with our own implementation of the standard mixed-model likelihood \citep{bates2015lme4}.
On the specification that the statsmodels \texttt{MixedLM} class can fit (without run slopes), our implementation reproduces the compound rates to within \EngineRateDiff\spl; \citet{seabold2010statsmodels} is the general reference for that package.
Intervals for compound rates are Wald intervals.
The interval for $\sigma_b$ is a profile-likelihood interval, cross-checked by a parametric bootstrap.
In 2026 each circuit is raced once, so $\sigma_b$ cannot distinguish a persistent circuit effect from race-to-race variation.
To separate them we also fit the 2022 to 2026 laps, with physical compounds throughout, with random tyre-age slopes for circuit and for event within circuit: the \PooledNRepeat{} circuits raced more than once share one circuit slope with standard deviation $\sigma_c$, and each race adds its own deviation with standard deviation $\sigma_e$.
The persistent share of the between-race slope variance is $\sigma_c^2/(\sigma_c^2+\sigma_e^2)$.
Both models are fitted by REML, and we test $\sigma_c=0$ (and, in 2026, $\sigma_b=0$) by a parametric bootstrap of the REML likelihood-ratio statistic (\CircuitNBoot{} replicates), because the null lies on the boundary of the parameter space and chi-square mixture approximations can be poor \citep{crainiceanu2004likelihood}.
The same bootstrap replicates give intervals for the components, their share and the difference between the circuit SD and the SD of the five compound rates.
Laps within a run are serially correlated and runs differ in wear, so the run terms make the effective sample size follow runs and circuits rather than laps.

We check calibration in two ways.
First, leaving each circuit out in turn, we ask whether the pooled model's predictive interval covers that circuit's own fitted rate for each compound.
Second, we simulate datasets from the fitted model and check whether intervals cover the true rates.
That check cannot reveal misspecification, so a separate recovery study simulates races with known rates from a process that Eq.~\eqref{eq:race} does not assume, with explicit fuel and track terms and pit stops that occur at random, earlier for slower cars, or when a run's own wear crosses a threshold (Appendix~\ref{app:recovery}).
The last scenario breaks the identifying assumption, because runs that degrade fast are cut short.
The run intercepts and slopes are also assumed independent of tyre age.
To check this we add each run's mean of its own tyre-age regressor as a fixed covariate, the device of Mundlak described by \citet[pp.~1056--1057]{bell2019fixed}, so that the slopes rest on variation within runs only, and test the added coefficients with a bootstrap over races or circuits (Appendix~\ref{app:mundlak}).
We compare three estimators of the same rate on the same laps.
The naive estimator regresses lap time on tyre age by compound with no controls and run-clustered errors.
The fuel-corrected estimator adds the assumed fuel effect and uses driver-run fixed effects, with no control for track evolution or traffic.
The third is Eq.~\eqref{eq:race}.

\subsection{Design analysis}\label{sec:design}

To learn what the data can detect, we simulate races from the fitted model of Eq.~\eqref{eq:race} with a known slope gap between two adjacent compounds, refit, and record how often a two-sided 5\% Wald test rejects.
We compare three allocations: balanced, with both compounds at every circuit; realistic, in which each run uses one of the three compounds Pirelli nominated for that 2026 race; and split, with the two compounds at disjoint circuits.
Runs last 20 laps and the variance components are the fitted ones, with sensitivity runs at the upper end of the circuit interval and with the run-slope standard deviation raised by half.
We compute the standard error analytically at known variances and check it with \PXSimReps{} simulated datasets per point that re-estimate the variances; the Type I error uses \PXSizeReps{} datasets at zero effect.

The split allocation has a floor.
If one compound runs only at $K_1$ circuits and the other only at $K_2$ different ones, each circuit adds its own slope deviation $b_e$ to every run there, so the contrast between the compounds has variance at least $\sigma_b^2(1/K_1+1/K_2)$ however many runs each circuit contributes.
More runs at the same circuits cannot remove this term; only more circuits can.
Within-race contrasts cancel $b_e$, and the realistic allocation allows them because compounds share races.
Because the simulated runs are idealised, we also report the minimum detectable effect implied by the fitted models' own standard errors, $(z_{0.975}+z_{0.80})\,\mathrm{SE}$, for each adjacent pair of compounds.

\subsection{Race and practice with one estimator}\label{sec:mgmt}

For each group $k$ (a season, event and compound label) we estimate a race slope and a practice slope with the estimator of Eq.~\eqref{eq:practice}, using the same code, tyre-age window (6 to 16 laps) and assumed fuel effect in both contexts.
Unlike Eq.~\eqref{eq:race}, which needs a common lap, this lets a race run be treated like a practice run.
Most practice race simulations start on used sets, whereas most race runs start on new ones, and used and new sets may degrade differently.
The primary design therefore keeps used sets only, in both contexts, as flagged in the timing feed.
We prefer this to our first design, which matched on the tyre age at which a run starts, because start age is a poor proxy for set state in races.
Of the \RaceRunsStartAge{} race runs that start at tyre age 5 or more, only \RaceRunsStartAgeUsed{} are used sets, and \RaceRunsStartAgeGap{} are new sets resuming after a gap in the lap record; in practice, \PracRunsStartAgeUsed{} of the \PracRunsStartAge{} runs that start at age 5 to 9 are used sets.
We report start-age matching, and the intersection of both rules, as robustness rows.
For a set of groups we pool laps within each context, so that the slope is $\sum xz/\sum x^2$ over all laps, and we report the gap $\Delta=\hat\beta^{\mathrm{race}}-\hat\beta^{\mathrm{prac}}$ and the ratio $\hat\beta^{\mathrm{race}}/\hat\beta^{\mathrm{prac}}$.
The gap is the primary quantity because the ratio of two noisy slopes is unstable and depends on the assumed fuel correction.
Intervals are 95\% percentile intervals from a bootstrap that resamples whole event-seasons (\NMatchedEventSeasons{} clusters in the primary design; 5,000 replicates) \citep{cameron2008bootstrap}.
A placebo applies the same design to FP2 against FP1 and FP3 together, where racing does not differ, to show how large a gap the design can produce by itself.
The placebo's orientation is arbitrary, so only its size matters.

Compound-specific management would show as a race--practice gap that grows with softness.
Within each event-season we estimate, in each context, the difference between the slopes of two compounds with lap fixed effects (ten-minute session bins in practice) and run fixed effects.
In a race the lap and run effects absorb the part of tyre age common to both compounds, so only the contrast is identified; two-way fixed effects remove lap- and run-level confounders only if effects are linear and additive \citep{imai2021twfe}.
The difference between the race and practice contrasts is a difference in differences.
We average it over events with inverse-variance weights that include an estimated between-event variance, take intervals from a bootstrap over events and $p$-values from a sign-flip permutation test with 10,000 draws, and adjust the three pairwise $p$-values by the Holm step-down method.
The primary test regresses the group-level race-minus-practice slope on softness (hard 0, medium 1, soft 2) with event-season fixed effects; it counts as evidence only if the bootstrap interval excludes zero and both one- and two-sided permutation $p$-values are below 0.05.
We obtain its minimum detectable effect by simulating the observed design.

The analyses differ in when they were specified.
The trend test and its decision rule were written into the analysis code before any slope was computed.
The decomposition into circuit and race components (Section~\ref{sec:mixed}) and the design analysis (Section~\ref{sec:design}) were added in response to an internal review, as were the recovery study and the check with run means of tyre age.
The recovery study exposed the shortcut described in Section~\ref{sec:within}, which our first analysis used; every result reported here uses the exact demeaning, Appendix~\ref{app:exact} compares the two, and the recovery criteria for the corrected estimator were set after the correction.
The used-set and start-age designs, the placebo and the difference in differences were added after an internal review of the first run; each was specified before it was run and none was tuned on its output.
The \MultiverseSpecs{} alternative specifications of Appendix~\ref{app:sens} and the 2023--24 telemetry extension of Appendix~\ref{app:tele} are exploratory and are reported because analysis choices can branch \citep{gelman2014crisis}.

\subsection{Forecasting race degradation from practice}\label{sec:forecast}

The forecast unit is an (event, compound) group.
For a weekend we compute the practice slope per compound and multiply it by a single scalar factor, the median ratio of race to practice slope over other weekends.
The target is the group's race slope, estimated by regressing lap time on tyre age and traffic within (event, lap) groups; a group needs at least three runs and a tyre-age spread of at least two laps across cars.
Groups this small cannot support the run-level random effects of Eq.~\eqref{eq:race}, so the target omits them and carries the pit-timing selection described in Section~\ref{sec:within}; the evaluation therefore measures how well practice predicts this per-race estimate (Section~\ref{sec:limits}).
In 2026 the factor for each weekend comes from the other weekends (leave-one-weekend-out).
For 2022 to 2025 we freeze the practice design and a factor of \FrozenFactor, both chosen on 2026 data, so these seasons are an out-of-season holdout (\HoldoutCells{} groups, \HoldoutWeekends{} weekends, label compounds) on which nothing was tuned.
The factor is smaller than the matched ratio of Section~\ref{sec:gap} because the two use different race estimators: here a cross-sectional fit within (event, lap) groups, there the within-run estimator on used sets.
Accuracy is the mean absolute error (MAE) in s/lap, with intervals from a bootstrap over weekends and paired differences between forecasts.
Baselines are the naive forecast (Sunday equals Friday), zero, the pooled race rate, a constant (the median race rate of other weekends) and the previous season's race rate at the same circuit.
Within 2026 (\TwentySixCells{} groups) we also select design and session weights inside each training fold, to measure how optimistic the tuned figure is.
